\section{Introduction}

\subsection{Document Purpose and Audience}

This document is the primary engineering specification for \EM{}, an iOS
application that reconstructs the analog photographic development chain as a
sequence of physically motivated computational transforms operating on
scene-referred RAW sensor data.

It is written to serve three audiences simultaneously, and each section is
marked for its intended reader where the distinction matters:

\begin{itemize}
\item \textbf{Implementing engineers} require the mathematics, the numerical
ranges, the GPU dispatch topology, the memory budget, and the persistence
schema. Sections~\ref{sec:colorspace} through~\ref{sec:schema} are normative for
implementation; where a formula appears with an equation number, that formula
is the contract.
\item \textbf{Technical reviewers and prospective investors} require the
argument that the approach is differentiated, defensible, and achievable.
Sections~\ref{sec:intro-thesis}, \ref{sec:background}, \ref{sec:roadmap},
\ref{sec:business}, and \ref{sec:risks} carry that argument and can be read
independently of the mathematics.
\item \textbf{Product and design stakeholders} require the interaction model
and the mapping from physical parameter to user-facing control.
Section~\ref{sec:ui} is the interface specification and
Table~\ref{tab:controlmap} is the authoritative translation table between
laboratory terminology and internal parameters.
\end{itemize}

\subsection{Scope}

In scope for the release covered by this document (v1.0, ``Still Photography''):
RAW capture on iPhone, the ten-stage development pipeline, non-destructive
editing with persisted parameter graphs, preset and recipe management, and
export to display-referred and intermediate formats.

Explicitly out of scope for v1.0, and deferred to the roadmap in
Section~\ref{sec:roadmap}: video capture and real-time video emulation, ACES
compliance certification, the macOS companion application, cloud
synchronization, user-authored film profiles, and any machine-learning
component. These are specified at architectural level only, sufficient to
guarantee that v1.0 decisions do not foreclose them.

\subsection{The Central Thesis}
\label{sec:intro-thesis}

The thesis of \EM{} is that the perceptual gap between mobile ``film camera''
applications and professional film emulation is not a gap of taste but a gap of
\emph{state space}. A lookup table is a memoryless, pointwise function
$f:\Rspace^3 \to \Rspace^3$. It cannot express any phenomenon whose output at a
pixel depends on the neighborhood of that pixel, on the local statistics of the
image, or on a physical quantity that was destroyed before the lookup table was
applied.

Nearly every phenomenon that a trained eye uses to identify photographic film
is exactly such a phenomenon:

\begin{itemize}
\item \textbf{Halation} is a spatial convolution whose strength depends on
scene luminance above a threshold and whose radius depends on wavelength. It is
not pointwise.
\item \textbf{Grain} is a stochastic process whose local variance depends on
local density and whose spatial correlation depends on emulsion structure. It
is not pointwise, and it is not deterministic.
\item \textbf{Interlayer inhibition} (the DIR coupler effect that gives color
negative film its characteristic edge behavior and saturation-versus-density
relationship) is a cross-channel spatial coupling. It is not pointwise.
\item \textbf{Highlight roll-off} depends on the linear scene-referred
radiance of the highlight, which a display-referred JPEG has already clipped.
The information required is not present in the input.
\end{itemize}

A pointwise operator applied to a display-referred image is therefore not an
approximation of film that happens to be imperfect; it is a class of function
that provably cannot represent the target. \EM{} is the argument that carrying
scene-referred data and modelling the actual transport is both feasible on
current mobile silicon and perceptually decisive.

\subsection{Design Principle: Parameter Honesty}

A secondary principle governs every user-facing decision in this document. We
call it \emph{parameter honesty}: no control is exposed unless it corresponds
to a physical quantity in the model, and no control is named for its perceptual
effect when a laboratory name exists.

This is not skeuomorphism. The application does not draw a fake developing
tank. It is the assertion that when a user increases \emph{development time},
the increase should propagate through the contrast index, which should steepen
the characteristic curve, which should raise mid-scale density, which should
change how the print stock's shoulder engages, and the resulting image should
differ from a naive contrast increase in precisely the ways that a real
push-processed negative differs. If the model is correct, the user learns
photography by using the application. If the model is a lookup table wearing
laboratory vocabulary, the user learns nothing and the vocabulary is a lie.

Concretely, parameter honesty implies three engineering constraints that
recur throughout this document:

\begin{enumerate}
\item \textbf{No orphan parameters.} Every element of the parameter vector
$\theta$ defined in Section~\ref{sec:apparch} appears in at least one equation
in Sections~\ref{sec:negative}--\ref{sec:aging}.
\item \textbf{Order is physical, not arbitrary.} Grain is applied to density,
not to the printed image, because grain is a property of the emulsion. Halation
is applied at exposure, not at output, because scattering happens before the
latent image forms. The pipeline order in Section~\ref{sec:arch} is derived,
not chosen.
\item \textbf{Defaults are measurements, not preferences.} The per-stock
parameters in Appendix~\ref{app:stockparams} are fit to published sensitometric
data and scanned reference wedges, not authored by eye.
\end{enumerate}

\subsection{Summary of Contributions}

The specific technical contributions specified in this document are:

\begin{enumerate}
\item A \emph{differentiable characteristic curve} formulation
(Section~\ref{sec:negative}) constructed from softplus primitives, giving
independent analytic control of toe sharpness, straight-line gamma, shoulder
sharpness, and $\Dmin$/$\Dmax$ asymptotes, with $C^\infty$ continuity and a
closed-form derivative required for both the tone-mapping Jacobian and the
grain variance model.
\item A \emph{filtered-Poisson grain model} (Section~\ref{sec:grain}) derived
from Campbell's theorem, which yields an analytic Wiener spectrum and therefore
permits exact spectral synthesis at any output resolution while preserving the
correct Selwyn granularity--density relationship and the correct
resolution-dependent appearance under downsampling.
\item A \emph{physically parameterized halation model}
(Section~\ref{sec:halation}) using an exponential point spread function whose
2D optical transfer function is available in closed form, together with a
third-order recursive filter approximation that evaluates in $O(1)$ per pixel
independent of radius.
\item A \emph{printer-light color control surface}
(Section~\ref{sec:printerlights}) calibrated in the industry-standard unit of
$0.025$ in $\log_{10}$ exposure per point, making \EM{} grades numerically
transferable to and from professional laboratory practice.
\item A \emph{render graph and precision policy} (Sections~\ref{sec:metal},
\ref{sec:rendergraph}) that sustains interactive editing on 48-megapixel
captures through resolution-decoupled preview, pointwise-chain LUT baking with
tetrahedral interpolation, and tile-memory kernel fusion.
\end{enumerate}

\subsection{Document Conventions}

Scene-referred linear quantities are written as radiometric exposures $E$ with
implied units of lux-seconds at the film plane. Logarithmic exposure is always
base ten, $x = \logE$, following sensitometric convention. Optical density is
$D = -\log_{10} T$ where $T$ is transmittance. Vectors over the three color
records are written in boldface, $\mathbf{D} = (D_R, D_G, D_B)^\top$.
Convolution is $*$, the two-dimensional Fourier transform is
$\mathcal{F}\{\cdot\}$, and spatial frequency in cycles per millimeter at the
film plane is $f$.

Requirement identifiers follow the pattern \reqid{FR-nn} for functional
requirements, \reqid{NFR-nn} for non-functional requirements, and \reqid{CR-nn}
for constraints, and are defined in Section~\ref{sec:requirements}.

Where a numerical value is given as a default, it is the shipping default and
is understood to be overridable by the stock profile.
